\chapter{Condition variables and semaphores}
\label{ch:condvars}

A mutex answers ``who may touch the data now?''. Threads also need to
\emph{wait for something}: for an element in a queue, for a thread to
finish, for the barber to wake up. Waiting correctly --- without burning
CPU and without missing the moment the condition becomes true --- is what
condition variables and semaphores are for. Your \code{pthread\_join} in
SWEB is one such wait; P1's elective asks for mutexes, condition variables
and semaphores in SWEB's user space.

\section{Waiting: spin or sleep}

\begin{lstlisting}
while (queue_empty(q)) {}          /* spinning: burns the CPU, and on one CPU
                                      the producer cannot even run */
\end{lstlisting}

Waiting should \emph{block}: take the thread off the CPU until another
thread says ``something changed''. Two problems must be solved together:
the waiter must not sleep through the signal (the \emph{lost wake-up}),
and the condition must be checked and the data accessed under a mutex.

\section{Condition variables}

\begin{lstlisting}
pthread_mutex_lock(&m);
while (!condition)                 /* ALWAYS a loop */
  pthread_cond_wait(&cv, &m);      /* atomically: unlock m and sleep;
                                      on wake-up: lock m again */
/* condition holds, and we hold m */
pthread_mutex_unlock(&m);

/* the other side */
pthread_mutex_lock(&m);
condition = true;
pthread_cond_signal(&cv);          /* wake one waiter (broadcast: all) */
pthread_mutex_unlock(&m);
\end{lstlisting}

A condition variable has no memory: a signal with nobody waiting is lost
--- by design. The \emph{condition} lives in ordinary shared variables,
protected by the mutex; the condition variable only provides ``sleep
until somebody might have changed it''.

\paragraph{Why \code{wait} must unlock atomically.} Suppose the waiter
unlocked, and only then went to sleep. In between, the signaller could
lock, set the condition, signal (nobody waits yet: lost) and unlock --- and
the waiter sleeps forever. \code{pthread\_cond\_wait} makes ``unlock and
sleep'' one step with respect to other threads. \module{05}'s
\code{lost\_wakeup.c} shows the broken and the fixed version.

\paragraph{Why \code{while}, not \code{if}.} Condition variables in
pthreads (and practically everywhere) have \emph{Mesa semantics}: a
signal only makes the waiter runnable; by the time it has re-acquired the
mutex, another thread may have consumed what it was woken for (a
\emph{stolen} wake-up). POSIX also allows \emph{spurious} wake-ups
without any signal. Rechecking in a loop handles both. (Under the
historical \emph{Hoare semantics} the signaller hands the mutex directly
to the woken thread, so the condition is guaranteed to hold --- simpler
to reason about, more expensive to implement.)

\paragraph{Signal or broadcast?} \code{signal} wakes one waiter: correct
when every waiter waits for the same thing and any one of them can
proceed. When waiters wait for different conditions on the same
variable, or when a change may allow several to proceed, use
\code{broadcast} --- a lost signal is a hang, an extra wake-up only costs
time.

\section{The bounded buffer}

The pattern to memorise (\module{05}, \code{bounded\_buffer.c}):

\begin{lstlisting}
void put(struct buf *b, int item) {
  pthread_mutex_lock(&b->m);
  while (b->count == CAP)
    pthread_cond_wait(&b->not_full, &b->m);
  b->item[(b->head + b->count++) % CAP] = item;
  pthread_cond_signal(&b->not_empty);
  pthread_mutex_unlock(&b->m);
}
int get(struct buf *b) {
  pthread_mutex_lock(&b->m);
  while (b->count == 0)
    pthread_cond_wait(&b->not_empty, &b->m);
  int item = b->item[b->head];
  b->head = (b->head + 1) % CAP;
  b->count--;
  pthread_cond_signal(&b->not_full);
  pthread_mutex_unlock(&b->m);
  return item;
}
\end{lstlisting}

\section{Semaphores}

Dijkstra's semaphore is a counter with two atomic operations:
\code{P} (\code{wait}, \code{sem\_wait}): wait until the value is
positive, then decrement; \code{V} (\code{signal}, \code{sem\_post}):
increment, waking a waiter if any. Unlike a condition variable, a
semaphore \emph{remembers}: a \code{post} with no waiter is kept in the
value. Three roles:

\begin{center}
\begin{tabular}{L{2.8cm}L{2cm}L{8.4cm}}
\toprule
\textbf{Role} & \textbf{Initial} & \textbf{Use} \\
\midrule
mutual exclusion & 1 & \code{wait}; critical section; \code{post} (a
``binary semaphore'' --- but anyone may \code{post}, unlike a mutex) \\
signalling & 0 & thread A \code{wait}s until thread B \code{post}s:
``B's event happened''; the order of the two calls does not matter \\
counting & $N$ & $N$ identical resources (buffer slots, connections) \\
\bottomrule
\end{tabular}
\end{center}

The bounded buffer with semaphores: \code{empty = CAP}, \code{full = 0},
\code{mutex = 1}; the producer does \code{wait(empty); wait(mutex); put;
post(mutex); post(full)}, the consumer the mirror image. The order of the
two \code{wait}s matters: \code{wait(mutex)} first and then blocking on
\code{empty} while holding it deadlocks the consumer.

A semaphore is easily built from a mutex and a condition variable
(\module{05}, part A) --- and conversely; they are equally powerful.

\section{Classic problems}

\paragraph{Readers--writers.} Many readers may read at once, a writer
needs exclusive access. The simple solution (readers count themselves;
the first reader locks writers out, the last one lets them in) lets a
steady stream of readers \emph{starve} writers. A writer-preferring lock
makes new readers wait as soon as a writer is waiting (\module{05},
part~B); a fair lock queues everybody in arrival order.

\paragraph{Sleeping barber.} A barber sleeps when there are no customers;
a customer wakes him, or waits in one of $N$ chairs, or leaves when all
chairs are taken. It combines signalling (customer wakes barber, barber
calls customer) with a counted resource (chairs) and needs care at the
edges: a customer arriving just as the barber decides to sleep must not
be lost, and shutdown must wake a sleeping barber (\module{05}, part~C).

\paragraph{Monitors.} A monitor bundles data, a mutex that every
operation holds, and condition variables --- the ``one object, one lock''
discipline as a language feature (Java's \code{synchronized} with
\code{wait}/\code{notify}). The bounded buffer above is a monitor written
by hand.

\begin{swebbox}
SWEB's \code{Condition} (\file{common/source/kernel/Condition.cpp}) is a
Mesa condition variable bound to a \code{Mutex}. \code{wait} puts the
thread on the waiters list \emph{while holding the list's lock}, releases
the mutex, and sleeps. The gap between unlocking the list and
\code{Scheduler::sleep()} is closed from the other side:
\code{Scheduler::wake} yields until the thread really is
\code{Sleeping} before it sets it \code{Running} --- no lost wake-up.
\code{Condition::wait} also asserts that the thread holds no other lock
while it sleeps. For \code{pthread\_join}: the joiner checks ``is the
thread finished?'' and waits on such a condition, under the lock that the
exiting thread takes to set ``finished'' and signal.
\end{swebbox}

\begin{keybox}
\begin{itemize}
  \item Wait by blocking; the condition lives in shared variables under a
    mutex.
  \item \code{cond\_wait} unlocks and sleeps atomically --- that is what
    prevents the lost wake-up.
  \item Always \code{while (!condition) wait}: Mesa semantics, stolen and
    spurious wake-ups.
  \item When unsure, \code{broadcast}.
  \item Semaphores remember posts: mutex (1), signal (0), count ($N$).
  \item Readers--writers: watch for writer starvation; barber: watch the
    edges and the shutdown.
\end{itemize}
\end{keybox}
